\section{Feature maps artifacts in ViT networks}

The following experiments aim to identify and characterize artifacts in feature maps
of both supervised and self-supervised ViT networks. According to the original paper, artifacts (or outlier) are defined as high-norm tokens appearing during inference primarily in low-informative back-
ground areas of images, that are repurposed for internal computations. We verify claims concerning the emergence and the information contents of these artifacts.

% Methodology
\subsection{Methodology}

This reproduction is implemented through original codebase that follows the procedure described in the paper and relies on public pretrained checkpoints loaded through \texttt{timm}, \texttt{open\_clip\_torch}, and Hugging Face \texttt{transformers}. Dataset loading is done through \texttt{torchvision} where possible otherwise is done through direct download from the dataset official website.

\subsubsection{Model descriptions}
All 4 experiments are conducted using ViT models without register as features extractor. Experiment 1 is run with \texttt{DeiT-III-B}, \texttt{DeiT-III ViT-B/16}, \texttt{DeiT-III ViT-L/16}, \texttt{OpenCLIP ViT-B/16}, \texttt{OpenCLIP ViT-L/14}, \texttt{DINO ViT-B/16}, and \texttt{DINOv2 ViT-g/14}. Experiment 2,3,4 employ \texttt{DINOv2 ViT-g/14} model only. Model details are in table \ref{tab:models}

\subsubsection{Datasets}
Datasets used in the following 4 experiments are \texttt{cifar10}, \texttt{cifar100}, \texttt{caltech101}, \texttt{dtd}, \texttt{aircraft}, \texttt{flowers102}, \texttt{food101}, and \texttt{pets}. See table \ref{tab:datasets} for details. 

\subsubsection{Hyperparameters}
The most important hyperparameter across these experiments is the minimum norm value a patch feature vector must have in order to be considered as an outlier patch. In the paper, this cutoff is chosen heuristically as 150; we use the same value unless otherwise specified.

\subsubsection{Experimental setup and code}

\textbf{Visualization of attention maps (Paper Figure 2)}. We run 6 ViT models on 4 images sampled from \texttt{Caltech101} and visualize the attention map from the last self-attention layer from each model. 

\textbf{DINO vs DINOv2 local feature norms (Paper Figure 3)}. We compare the local features norms distribution for \texttt{DINO ViT-B/16} and \texttt{DINOv2 ViT-g/14} both on a single image and on a batch of 1000 images sampled from \texttt{Caltech101}.

\textbf{Local properties of outlier patches (Paper Figure 5)}. We compare the distribution of cosine similarity between input patches embeddings and their 4 neighbors for both normal patches and outlier patches. In addition we perform linear probing for local information on both normal and outlier patches by training:
namely (i) a linear classifier for patch-position prediction and (ii) a linear regressor for patch-pixel reconstruction.
Probings are done with 500 images sampled from \texttt{Caltech101}. Normal and outlier patches are both present in the training split while we test them separately on evaluation. We report the hyperparameter of the 2 linear models: 
\begin{table}[h]
    \centering
    \label{tab:figure5b-hparams}
    \resizebox{\textwidth}{!}{%
    \begin{tabular}{llllllllllll}
        % \toprule
        Train fraction & Max epochs & Batch size & Learning rate & Optimizer & Weight decay & Early stopping patience & Position probe loss & Reconstruction probe loss & Feature dimension & Number of positions & Pixel target dimension \\
        \midrule
        0.75 & 80 & 512 & 0.001 & AdamW & $10^{-4}$ & 3 & Cross-entropy & MSE & 1536 & 256 & 588 \\
        \bottomrule
    \end{tabular}%
    }
\end{table}

The model used as features extractor is \texttt{DINOv2 ViT-g/14}

\textbf{Global properties of outlier patches (Paper Table 1)}
We compare the class predictive capacity of different patch tokens: for each image we extract three candidate representations: the \texttt{[CLS]} token, one normal patch token, and one outlier patch token. We then trains three separate multinomial logistic regression classifiers, one per representation type. Among the 14 datasets chosen in the paper we pick 8 that are easier to download and store locally. We fit the logistic model on these 8 datasets.
For each dataset we sample 500 images for training and 200 images for evaluation. If a certain image doesn't contain any outlier patch (norma > 150) we pick the the patch with the highest norm as outlier. The model used as features extractor is \texttt{DINOv2 ViT-g/14}


\subsubsection{Computational requirements}
The experiments executed locally on an Apple M4 Pro machine with 24\,GB unified memory. The mps backend is used whenever is possible; when some torch implementation are known to be unstable we fall back to cpu.
For experiment 1 the runtime is neglectable, for experiment 2 is $\approx 20 min$, for experiment 3 $\approx 15 min$, for experiment 4 $\approx 1h $.  



% Results
\subsection{Results}
\label{sec:results}
The 4 experiments are intended to verify 4 qualitative claims of the original paper in relation to the main characteristic and information stored inside outlier patches tokens representation. 

\subsubsection{Results reproducing original paper}

\textbf{Claim 1: all models but DINO exhibit peaky outlier values
in the attention maps.} Our attention maps reproduce faithfully those in the paper where all models but DINO exhibit peaky outlier values (see \ref{fig:outlier_maps})

\textbf{Claim 2: DINOv2 has a few outlier patches, whereas DINO does not present these artifacts.}. Across 1000 images our experiment confirms that DINO does not present outlier patches whereas DINOv2 has them. The norm distribution of features extracted by DINO and DINOv2 resembles the distribution presented in the paper (see \ref{fig:dino_vs_dinov2})

\textbf{Claim 3: high-norm tokens appear where patch information is redundant and hold little local information}. Our experiment confirms that high norm tokens appears on patches that are very similar to their neighbors (see histogram in fig \ref{fig:neighbor_cosine}). Local probe results show that these high norm tokens hold little local information compared to normal norm tokens (see table in fig \ref{fig:neighbor_cosine}) since they achieve lower scores both in patch position prediction and patch reconstruction. 

\textbf{Claim 4: high-norm tokens hold global information}
Our experiment confirms that outlier tokens have higher accuracy on global class prediction, compared to normal tokens, across the 8 datasets that we tasted. They contains more global information about the global image. (see table \ref{tab:global-probes}). 

\begin{table}[h]
    \centering
    \footnotesize
    \caption{Top-1 accuracy (\%) for classification from the \texttt{[CLS]} token, one normal patch token, and one outlier patch token.}
    \label{tab:global-probes}
    \begin{tabular}{lcccccccc}
        \toprule
        Token & CIFAR10 & CIFAR100 & Caltech101 & DTD & FGVCAircraft & Flowers102 & Food101 & OxfordIIITPet \\
        \midrule
        \texttt{[CLS]} & 99.0 & 85.0 & 92.0 & 77.0 & 51.5 & 99.5 & 79.0 & 92.5 \\
        Normal patch   & 82.5 & 29.5 & 41.0 & 44.0 & 3.5  & 23.5 & 10.5 & 14.5 \\
        Outlier patch  & 99.0 & 83.0 & 94.5 & 79.0 & 45.0 & 99.0 & 73.5 & 84.5 \\
        \bottomrule
    \end{tabular}
\end{table}




% Discussion
\subsection{Discussion}

All our experiments confirm claims from the paper regarding the definition of outlier patches and their information content compared to normal patches. 

\subsubsection{What was easy}
The paper describes clearly the pipeline followed to produce the results claimed even if the code was not released. 

\subsubsection{What was difficult}
For the first 3 experiments the paper doesn't specify explicitly the datasets used but nevertheless results hold with the dataset chosen by us. Regarding the last experiment on different datasets, it was impractical for us to run probes on the full size of each of datasets. Fortunately, since ViT models often produce highly informative representations, probing on a portion of each dataset is enough to verify qualitatively the global properties of outlier patches.  

